% ============================================================
% CAPITOLO 29 - SPECIE RARE (FACOLTATIVO)
% ============================================================
\chapter{Specie rare (facoltativo)}
\label{chap:specie_rare}

Le ultime due specie completano il catalogo delle quadriadi, ma non trovano applicazione nell'armonia diatonica tradizionale del basso dato. Il capitolo si può saltare senza perdere il filo.

\section{Sesta specie: minore con settima maggiore}
\label{sec:sesta_specie_dettaglio}

La sesta specie è formata da una \textbf{triade minore} e una \textbf{settima maggiore}. È un accordo \textbf{secondario e artificiale}, di ambito moderno: nell'armonia diatonica classica era considerato un accordo alterato.

\begin{regola}[Grado e risoluzione della sesta specie]
	\label{reg:sesta_specie}
	La sesta specie si trova sul \textbf{I grado della scala minore armonica e melodica}, con le sigle $m^{maj7}$ e $m\Delta$ (per esempio La-Do-Mi-Sol$\sharp$). Al basso ha la fondamentale e risolve \textbf{di norma sul IV grado}.
\end{regola}

La settima è la sensibile della tonalità: non può scendere, quindi \textbf{sale} alla tonica (Sol$\sharp$ $\rightarrow$ La) e non si raddoppia. È l'unica specie in cui la settima sale. Va preparata come le altre settime secondarie.

\begin{esempio}[Sesta specie in La minore: \sigla{I}{7}{} $\rightarrow$ IV]
	\label{ex:sesta_specie}
	Nello stato fondamentale e nei rivolti la settima Sol$\sharp$ sale a La, che è la quinta del IV (Re-Fa-La). Nel terzo rivolto il basso Sol$\sharp$ sale a La e il IV è in secondo rivolto.
	\musica{es_cap13_minmagg_la_1}
	\musica{es_cap13_minmagg_la_2}
\end{esempio}

\section{Settima specie: aumentata con settima maggiore}
\label{sec:settima_specie_dettaglio}

La settima specie è formata da una \textbf{triade aumentata} e una \textbf{settima maggiore}. È anch'essa \textbf{secondaria e artificiale}, di ambito moderno o jazzistico, e considerata alterata nella prassi classica.

\begin{regola}[Grado e risoluzione della settima specie]
	\label{reg:settima_specie}
	La settima specie si trova sul \textbf{III grado della scala minore armonica e melodica}, con le sigle $maj^7(\sharp5)$ e $+\Delta$. Al basso ha la fondamentale e risolve \textbf{di norma sul VI grado}.
\end{regola}

Due note hanno una tendenza: la settima Si \textbf{scende} a La, la quinta eccedente Sol$\sharp$ \textbf{sale} a La. Nessuna delle due si raddoppia.

\begin{esempio}[Settima specie in La minore: \sigla{III}{7}{} $\rightarrow$ VI]
	\label{ex:settima_specie}
	Nel primo rivolto (basso Mi) il basso sale a Fa; nel secondo (basso Sol$\sharp$) sale a La e il VI è in primo rivolto; nel terzo (basso Si) scende a La.
	\musica{es_cap13_ecc_la_1}
	\musica{es_cap13_ecc_la_2}
\end{esempio}
